\scanpage{45}
se, if the index theorem is proved for $X$, then $A_{\chi}(X,\inserted{\unsure{\xi}})=B_{\chi}(X)$ is again independant of the polarization, and analoguous remark for $C_{\chi}(X,\inserted{\unsure{\xi}})$.

When speaking about condition $C_{\chi}(X,\inserted{\unsure{\xi}})$, emphasize at once its stability properties by products (the proof I suggested works indeed) specialization (with possible change of characteristics), hyperplane or more generally linear sections. Give an extra proposition for the relations with the property $A$, via a formal proposition as follows:

\uline{Proposition}\quad Equivalent conditions on $X$(\uline{polarized} variety):

(i) $C_{\chi}(X)$

(ii) $C_{\chi}(Y)$ and $A_{\chi}(X\mathrm{x}X)^{\deleted{\unsure{0}}}$, \deleted{where the exponent $^0$}

(ii bis) $C_{\chi}(Y)$ and $A_{\chi}(X\mathrm{x}X)^0$, where the exponent $^0$ means that one restricts oneself to expressing condition $A$ for the homomorphism in critical dimension
\[
H^{2n-2}\longrightarrow H^{2n+2}.
\]

(iii) $C_{\chi}(Y)$ and $A_{\chi}(X\mathrm{x}Y)$.

(iii bis) $C_{\chi}(Y)$ and $A_{\chi}(X\mathrm{x}Y)^0$, where the exponent $^0$ means that one restricts oneself to expressing condition $A$ in critical dimension $H^{(2n-1)-1}\longrightarrow H^{(2n-1)+1}$.

(iv) \deleted{For every $\unsure{n|O|0}$} $C_{\chi}(Y)$, and for every $i\leq n-1$, the natural homomorphism $H^i(Y)\mathbin{-}H^i(X)$ left inverse to $\varphi^i:H^i(X)\mathbin{-}H^i(Y)$ (induced by $\Lambda_X\varphi_{*}$) is induced by an algebraic correspondence class (inducing whatever it likes on the other $H^j(Y)$).

(iv bis) $C_{\chi}(Y)$, and for $j\geq n+1$, the natural homomorphism $H^j(X)\longrightarrow H^{j\inserted{-2}}(Y$ right inverse to $\varphi_{j-2}:H^{j-2}(Y)\mathbin{-}H^j(X)$ \deleted{is} (induced by $\varphi^*\Lambda_X$) is induced by an algebraic correspondence class (inducing whatever it likes on the other $H^i(X)$).

\uline{Corollary}\quad \unsure{These|These} conditions are also equivalent to
\[
\text{(v)}\qquad A_{\chi}(X\mathrm{x}X)^0+A_{\ell}(Y\mathrm{x}Y)^0+A_{\ell}(Z\mathrm{x}Z)^0+\qquad,
\]
where $X\supset Y\supset Z$ is a decreasing sequence of \inserted{nonsingular} hyperplane sections
\mspage{44}
